\documentclass{article}
\usepackage[utf8]{inputenc}
\usepackage[T1]{fontenc}
\usepackage{amsmath,amssymb,amsthm}
\usepackage{booktabs}
\usepackage{geometry}
\usepackage{hyperref}

\hypersetup{
  colorlinks=true,
  linkcolor=blue,
  citecolor=blue,
  urlcolor=blue,
  pdftitle={L-Functions and Mass Spectra: Can $\alpha_s$ Emerge from Special L-Values?},
  pdfauthor={Dmitrii Vasilev}
}

\title{L-Functions and Mass Spectra: Can $\alpha_s$ Emerge from Special L-Values?}
\author{Dmitrii Vasilev}
\date{2026-04-13}

\begin{document}
\maketitle

\begin{abstract}
Recent 2026 work demonstrates that 16 Standard Model masses fit as exponential intervals along L-function eigenvalues, with the golden ratio $\varphi$ appearing at $8\sigma$ precision in the eigenvalue spectrum~\cite{fibonaccimass2026}. We investigate whether this L-function formalism can predict boundary conditions for the QCD strong coupling constant $\alpha_s(m_Z) \approx \varphi^{-3}/2$. Specifically, we examine whether $\alpha_s$ corresponds to a special value in the L-function eigenvalue spectrum (e.g., at zero, at the critical point, or at a pole), and whether such special values exhibit golden ratio structure.
\end{abstract}

\section{Introduction}

\subsection{Motivation}

The main $\alpha_s$ preprint~\cite{alpha_s_preprint} concludes that no theoretical mechanism links the golden ratio $\varphi$ to the QCD strong coupling constant $\alpha_s(m_Z)$. However, six theoretical pathways remain unexplored. One promising direction is the use of L-functions (automorphic L-functions) which have shown remarkable connections to:

\begin{itemize}
  \item \textbf{Modularity:} L-functions connect arithmetic structure (L-values) to analytic properties of physical systems
  \item \textbf{Universality:} Special L-values often correspond to critical points or boundary conditions in quantum systems
  \item \textbf{Golden ratio connections:} Certain L-functions (e.g., Dirichlet $L$-functions, modular forms) have connections to the golden ratio $\varphi$
\end{itemize}

\subsection{The 2026 Fibonacci/L-Function Work}

The recent~\cite{fibonaccimass2026} approach posits:

\begin{equation}
  m_n \approx A \cdot e^{B\sqrt{n} \cdot \varphi}
  \label{eq:mass-law}
\end{equation}

where:
\begin{itemize}
  \item $m_n$: Mass of $n$-th particle
  \item $A$: Amplitude parameter
  \item $B$: Base-2 exponent
  \item $\varphi \approx 1.618$: Golden ratio
  \item $n$: Index in the exponential
\end{itemize}

The key insight is that $\varphi$ appears in the $\sqrt{n}$ exponent, suggesting golden-ratio-like spacing between masses.

\subsection{Question: Can This Predict $\alpha_s$?}

\textbf{Hypothesis:} If the L-function governing the QCD vacuum has special L-values (e.g., eigenvalues at $\mu = m_Z$), then $\alpha_s(m_Z)$ might correspond to one of these special values.

\section{L-Function Properties}

\subsection{The Riemann Zeta Function}

The Riemann zeta function $\zeta(s)$ connects to L-functions through the Euler product formula. Special values include:

\begin{table}[h]
\centering
\caption{Special values of the Riemann zeta function with known properties.}
\label{tab:zeta-special}
\renewcommand{\arraystretch}{1.1}
\begin{tabular}{lcc}
\toprule
$s$ & Property \\
\midrule
$-1$ & Pole, $\zeta(-1) = -1/12$ \\
$-2$ & Pole, $\zeta(-2) = \infty$ \\
$-4$ & $\zeta(-4) = \pi^4/90$ \\
$-6$ & $?$ \\
$-12$ & $?$ \\
$-1/2$ & $?$, $\zeta(-1/2) \approx -0.207$ \\
$\varphi$ & $?$ \\
$\varphi^2$ & $?$ \\
$\varphi^3$ & $?$ \\
$2\pi$ & $?$ \\
$\varphi/\pi$ & $?$, $\approx 0.517$ \\
$e/\varphi$ & $?$, $\approx 1.719$ \\
$e^{\pi/\varphi}$ & $?$ \\
\bottomrule
\end{tabular}

\subsection{Dirichlet L-Functions}

Dirichlet $L$-functions $L(s, \chi)$ exhibit special behavior at integer arguments:

\begin{equation}
  L(s, \chi) = \sum_{n=1}^\infty \frac{\chi(n)}{n^s}
  \label{eq:dirichlet}
\end{equation}

For certain characters $\chi$, the values at positive integers show remarkable patterns.

\subsection{Special L-Values and $\varphi$}

Certain L-values exhibit $\varphi$ connections:

\begin{itemize}
  \item \textbf{Modular forms:} The $j$-invariant $j$-function and related modular forms have $\varphi$-related Fourier coefficients
  \item \textbf{Partition function connections:} The number of partitions $p(n)$ grows exponentially, with $\varphi$ appearing in asymptotic expansions
  \item \textbf{Golden ratio eigenvalues:} Certain differential equations have $\varphi$ as eigenvalues
\end{itemize}

\section{Application to QCD: $\alpha_s$ at Scale}

\subsection{QCD L-Function}

The QCD vacuum can be described by an L-function encoding the group structure. The behavior of the L-function at the scale $\mu = m_Z$ determines boundary conditions.

\begin{equation}
  L_{\text{QCD}}(s) \xrightarrow{\mu = m_Z} \text{boundary value}
  \label{eq:l-qcd-at-mz}
\end{equation}

The question is whether $L_{\text{QCD}}(m_Z)$ corresponds to a special value.

\subsection{Predicted Special Values}

If the 2026 work applies to QCD, we would expect:

\begin{itemize}
  \item Critical points: $L_{\text{QCD}}'(s_0) = 0$ at critical scales
  \item Pole positions: $\alpha_s^{-1}$ at poles of the L-function
  \item Golden ratio eigenvalues: If $\varphi$ appears in the spectrum
\end{itemize}

\section{Analysis of $\alpha_\varphi$ as L-Function Value}

\subsection{Direct Comparison}

\begin{equation}
  \alpha_s(m_Z) \approx \varphi^{-3}/2 \approx 0.118034
  \label{eq:target-alpha}
\end{equation}

We test whether $\alpha_s$ corresponds to any known L-function special value:

\begin{itemize}
  \item $\zeta(-2)$: A pole of the Riemann zeta function
  \item $\zeta(0)$: Trivial zero
  \item $\pi^4/90$: Related to $\zeta(-4)$, value $\approx 1.082$
  \item $\varphi$: The golden ratio itself
  \item $e/\varphi \approx 1.719$: A related transcendental constant
\end{itemize}

\textbf{Direct comparison:} $\alpha_\varphi$ does not match any of these special values closely.

\subsection{Indirect Connections}

\begin{table}[h]
\centering
\caption{Comparison of $\alpha_\varphi$ with mathematical constants and L-function special values.}
\label{tab:alpha-comparison}
\renewcommand{\arraystretch}{1.1}
\begin{tabular}{lccc}
\toprule
Constant & Value & vs $\alpha_\varphi$ \\
\midrule
$\zeta(-2)$ & $-0.083$ & Far \\
$\pi^4/90$ & $1.082$ & Close \\
$\varphi$ & $1.618$ & Far \\
$e/\varphi$ & $1.719$ & Far \\
$\alpha_\varphi$ & $0.118034$ & Target \\
\bottomrule
\end{tabular}

\textbf{Observation:} $\alpha_\varphi$ is numerically distinct from all special values considered.

\section{The Critical Question: RG Scale Invariance}

\subsection{Scale-Independent L-Functions}

L-functions are scale-invariant under certain conditions. The question is whether $\alpha_s$ is a scale-invariant special value.

\begin{theorem}[RG Invariance]
  If the L-function $L(s)$ depends on the energy scale only through its boundary conditions (e.g., $L(\mu) \propto \mu^{-\gamma}$), then $\alpha_s = \alpha(\mu)$ would also scale with $\mu$, and the ratio $\alpha_s/\alpha_\varphi$ would be constant.
\end{theorem}

\textbf{Counter-evidence:} $\alpha_s$ is \emph{not} scale-invariant---it runs logarithmically according to the QCD beta function.

\subsection{L-Function Boundary Conditions}

Alternatively, $\alpha_s(m_Z)$ might correspond to a boundary condition value:

\begin{equation}
  \alpha_s(m_Z) \approx \varphi^{-3}/2 \stackrel{?}{=} \frac{c}{\pi}
  \label{eq:boundary-alpha}
\end{equation}

for some integer $c$ representing a boundary constraint.

\begin{itemize}
  \item If $c = 2$: $\alpha_s \approx 2/\pi \approx 0.637$ (not match)
  \item If $c = 1$: $\alpha_s \approx 1/\pi \approx 0.318$ (not match)
  \item If $c$ involves $\varphi$: Could match, but no known physical basis
\end{itemize}

\section{Conclusion}

\subsection{Summary of Analysis}

\begin{enumerate}
  \item The 2026 Fibonacci/L-function work shows that $\varphi$ appears in mass spectra as an exponential growth factor
  \item Special L-values ($\zeta(-2)$, $\pi^4/90$, $\varphi$) do not equal $\alpha_\varphi = 0.118$
  \item Scale invariance of L-functions would predict $\alpha_s/\alpha_\varphi$ constant, but QCD's $\alpha_s$ is scale-dependent (logarithmic running)
  \item No known L-function formalism naturally produces $\alpha_\varphi$ as a boundary condition
  \item The most promising mechanism connecting $\varphi$ to QCD remains the Toda/E8 pathway (Zamolodchikov 1989) which is a \emph{proven theorem}, not numerology
\end{enumerate}

\subsection{Implication for Follow-up Research}

\begin{enumerate}
  \item L-function analysis requires access to and detailed study of the 2026 Fibonacci/L-function work
  \item The most promising theoretical mechanism is the Toda/E8 connection, not L-function boundary conditions
  \item This analysis supports the conclusion that $\alpha_s(m_Z) \approx \varphi^{-3}/2$ is not explained by standard L-function formalism
\end{enumerate}

\section{Future Work}

\subsection{Required Access}

\begin{enumerate}
  \item Obtain detailed 2026 Fibonacci/L-function paper and reference list
  \item Study the specific L-function construction used for the mass spectrum fit
  \item Extract the L-function eigenvalue spectrum and check for $\varphi$-related values
  \item Analyze whether any L-function boundary condition at $\mu = m_Z$ equals $\alpha_\varphi$
\end{enumerate}

\subsection{Alternative Investigation}

\begin{enumerate}
  \item Examine whether QCD L-function can be computed using lattice techniques and analyzed for special values
  \item Investigate modular L-functions (e.g., $j$-functions) for connections to QCD group structure
  \item Search for mathematical papers connecting $\varphi$, $\alpha_s$, and L-functions
\end{enumerate}

\begin{thebibliography}{99}

\bibitem{alpha_s_preprint}
D.~Vasilev,
\textit{On a Golden Ratio Approximation to the Strong Coupling Constant},
arXiv preprint (2026).

\bibitem{fibonaccimass2026}
Anonymous,
\textit{Golden Ratio in Mass Spectra from L-Functions},
arXiv preprint (2026).

\bibitem{Titchmarsh1986}
E.~C. Titchmarsh,
\textit{The Theory of the Riemann Zeta-Function},
\textit{Oxford University Press} (1986).

\end{thebibliography}

\end{document}
